\section{Motivation and the question being tested}
The starting picture is a three-dimensional universe represented by the
three-dimensional boundary of a four-dimensional ball. Each layer is a
complete spatial slice at one time. Successively larger layers represent
later stages of expansion. A familiar two-dimensional balloon surface is
a lower-dimensional analogy for this construction.

The important observational object is the \emph{view line}: the image of
a received light ray, traced backward from the observer. Its earlier events
lie on earlier spatial layers. A diagram therefore separates the universe
at a specified time from the collection of emission events visible to an
observer. The physical basis is standard null propagation and cosmological
distance geometry \citep{Hogg1999,DavisLineweaver2004}.

The motivation for a local-mass extension is equally concrete. Can the same
representation show spatial distortion, gravitational time dilation, and
the changes in view lines without assigning conflicting meanings to drawing
radius? We answer this for a controlled linear example. We first specify a
spacetime metric and a consistent matter perturbation. We then derive the
clock and ray equations, and finally choose drawings of those quantities.

This is a reorganization and extension of the v0.4 manuscript. Its core
distinction between light-cone observation and simultaneous separation is
retained. The original text and three figures are preserved in
\texttt{archive/v0.4/}. The earlier proposed Bianchi/vorticity program is
deferred while the local-gravity question is developed. Neither version
establishes a priority claim for polar cosmological diagrams.

\paragraph{Scope.}
The background here is positively curved, simply connected, and expanding
over the plotted interval. It contains pressureless matter and no cosmological
constant. This choice permits explicit calculations; it is not a fit to the
observed universe. The local source is a small, evolving dust overdensity
with compensating surroundings. It is not a star of constant mass, a point
singularity, or a black hole. The distinction is essential when interpreting
the resulting geometry.
